\section{Experiments} \label{sec:experiments}

We conduct four experiments of increasing difficulty to evaluate how
defense strategies perform as attack conditions become less predictable.
All experiments use the same six MAVLink injection attacks
(Table~\ref{tab:attacks}), the same platform simulator at 10~msg/s, and
a 10-second baseline training phase.  Randomized experiments use seed~42
for reproducibility.

\subsection{Experiment Summary}

Table~\ref{tab:results} summarizes the key metrics across all four
experiments.

\begin{table}[t]
\centering
\caption{Results across the four experimental conditions.}
\label{tab:results}
\begin{tabular}{lcccc}
\toprule
\textbf{Metric} & \textbf{E1} & \textbf{E2} & \textbf{E3} & \textbf{E4} \\
\midrule
Defense mode     & None & Markov & Markov & IF Adaptive \\
Attack schedule  & Static & Static & Random & Random \\
Block rate       & 0\%  & 40.3\% & 43.1\% & \textbf{62.5\%} \\
Messages blocked & 0    & 1{,}398  & 2{,}286  & 2{,}279 \\
Messages passed  & --   & 2{,}073  & 3{,}020  & 1{,}370 \\
Detection latency (s) & 41.2 & 60.5 & 3.6 & 5.2 \\
Total alerts     & 1    & 1    & 1    & 1 \\
\bottomrule
\end{tabular}
\end{table}

% -----------------------------------------------------------------------
\subsection{E1: Static Attacks, No Defense}

\textbf{Setup.}  Six attacks run sequentially at 200~msg/s for 10
seconds each, in fixed order.  The detector monitors and logs but does
not block any traffic.

\textbf{Command.}
\begin{verbatim}
python run_demo.py --mode all \
    --defense-mode none --no-live-plot
\end{verbatim}

\textbf{Results.}  Figure~\ref{fig:e1} shows the three-panel detection
report.  The moving-average accuracy drops sharply when each attack
begins and partially recovers between attacks as normal simulator traffic
resumes.  The anomaly score drops during attack periods, confirming that
the Isolation Forest correctly identifies the injected traffic as
anomalous.  No messages are blocked because no defense is active.

The detection latency of 41.2~seconds reflects the time from the end of
baseline training to the first alert; the detector correctly identifies
anomalous conditions but has no mechanism to act on them.

\begin{figure}[t]
\centering
\includegraphics[width=\columnwidth]{out/E1/accuracy_plot.png}
\caption{E1: Static attacks with no defense.  Accuracy drops during each
attack episode and partially recovers between episodes.}
\label{fig:e1}
\end{figure}

% -----------------------------------------------------------------------
\subsection{E2: Static Attacks, Markov Defense}

\textbf{Setup.}  Same static attack schedule as E1, but with the Markov
transition-probability defense enabled (threshold = 0.05).

\textbf{Command.}
\begin{verbatim}
python run_demo.py --mode all \
    --defense-mode markov \
    --defense-threshold 0.05 --no-live-plot
\end{verbatim}

\textbf{Results.}  Figure~\ref{fig:e2} shows that the Markov defense
blocks 40.3\% of attack traffic overall.  However, the blocked messages
are almost entirely concentrated on one attack type: ping floods
(message type~4) from source~251.  The defense fails to block heartbeat
floods, MITM identity spoofs, replay patterns, and command injections
because these attacks either use message types that have legitimate
Markov transitions or introduce new source IDs that the transition model
does not track.

This demonstrates the limitation of a pattern-based defense that operates
on a single feature (message-type ordering): it catches attacks that
disrupt the learned transition sequence but misses attacks that
manipulate other traffic dimensions.

\begin{figure}[t]
\centering
\includegraphics[width=\columnwidth]{out/E2/accuracy_plot.png}
\caption{E2: Static attacks with Markov defense.  Partial blocking of
ping floods, but most attack types pass through.}
\label{fig:e2}
\end{figure}

% -----------------------------------------------------------------------
\subsection{E3: Randomized Attacks, Markov Defense}

\textbf{Setup.}  Attacks are randomized: shuffled order, each episode
split into 5--13 segments with independently sampled rates (60--300+
msg/s) and durations, with random idle gaps between episodes.  The
Markov defense remains active.

\textbf{Command.}
\begin{verbatim}
python run_demo.py --mode all \
    --defense-mode markov \
    --defense-threshold 0.05 \
    --randomize-attacks --random-seed 42 \
    --no-live-plot
\end{verbatim}

\textbf{Results.}  Figure~\ref{fig:e3} shows a more volatile accuracy
trace.  The Markov defense achieves a slightly higher overall block rate
(43.1\%) because the randomized schedule happens to produce more of the
transitions it can catch.  However, it still blocks only a subset of
attack types (ping and param-request floods) and misses identity spoofs,
replay attacks, and command injections.

The accuracy trace is visibly more irregular than E2, reflecting the
non-stationary attack traffic.  Critically, accuracy does not recover to
baseline levels between attack episodes because the Markov defense cannot
prevent the majority of attack traffic from reaching the detector.

\begin{figure}[t]
\centering
\includegraphics[width=\columnwidth]{out/E3/accuracy_plot.png}
\caption{E3: Randomized attacks with Markov defense.  The accuracy trace
is volatile and the defense blocks only a subset of attack types.}
\label{fig:e3}
\end{figure}

% -----------------------------------------------------------------------
\subsection{E4: Randomized Attacks, Isolation Forest Adaptive Defense}

\textbf{Setup.}  Same randomized attack schedule as E3 (seed~42), but
with the Isolation Forest adaptive defense replacing the Markov defense
(anomaly-score threshold = 0.3).

\textbf{Command.}
\begin{verbatim}
python run_demo.py --mode all \
    --defense-mode adaptive \
    --randomize-attacks --random-seed 42 \
    --no-live-plot
\end{verbatim}

\textbf{Results.}  Figure~\ref{fig:e4} shows a substantial improvement.
The Isolation Forest adaptive defense achieves a 62.5\% block rate,
blocking traffic across \emph{all} attack types and \emph{all} source
systems.  Unlike the Markov defense, which only caught transitions it
had learned, the Isolation Forest operates on the full 11-dimensional
feature space and detects deviations in rate, source count, entropy,
inter-arrival timing, and per-type ratios simultaneously.

The accuracy trace shows sharper recovery after attack episodes because
anomalous messages are blocked before they corrupt the Markov prediction
window.  The cumulative blocked count increases across all attack
types, confirming that the defense generalizes to identity spoofs,
command injections, and replay patterns---attack vectors that the Markov
defense missed entirely.

\begin{figure}[t]
\centering
\includegraphics[width=\columnwidth]{out/E4/accuracy_plot.png}
\caption{E4: Randomized attacks with Isolation Forest adaptive defense.
Higher block rate, broader coverage across attack types.}
\label{fig:e4}
\end{figure}

% -----------------------------------------------------------------------
\subsection{Discussion}

The progression from E1 to E4 demonstrates two key findings:

\textbf{1. Pattern-based defenses are brittle.}  The Markov defense
blocks only attacks that disrupt the learned message-ordering pattern.
Attacks that inject traffic using legitimate message types at plausible
rates (e.g., identity spoofing, command injection) evade it entirely.

\textbf{2. Feature-based ML defenses generalize.}  The Isolation Forest,
trained on an 11-dimensional feature vector capturing rate, entropy,
source diversity, and transition statistics, detects and blocks attack
traffic regardless of which specific parameter the attacker varies.
Under the same randomized attack schedule, the Isolation Forest achieves
a 62.5\% block rate compared to the Markov defense's 43.1\%, with
coverage across all six attack types.

These results validate the harness design: by running different defenses
against the same attack schedule under identical conditions, we can
make direct, fair comparisons that would not be possible with ad-hoc
evaluation setups.
